\section{Integral tame-character boundary lifting at the added split prime}
\label{sec:cubic-hsieh-tame-type}

The coefficient prime is $p>3$, and $q\ne p$ is the added
split rational prime of AC-G67-measure.  Write
$\vartheta_q:\mathbb Z_q^\times\to\mathcal O'{}^\times$
for its tame character of order $p^r$.  All the original
Hsieh and repaired nonvanishing inputs in that proposition
remain in force.  At $q$ use
$U(V)(\mathbb Q_q)=\mathrm{GL}_3(\mathbb Q_q)$, with upper
Borel $B$, upper unipotent group $N$, diagonal torus $T$, and
longest Weyl element $w$.  The opposite Iwahori $I^-$ consists
of matrices whose reduction modulo $q$ is lower triangular.
Its character is
\[
 \lambda(\operatorname{diag}(t_1,t_2,t_3))
       =\vartheta_q(t_1)^{-1}\vartheta_q(t_3).
                                                        \tag{G67.2}
\]
For the similitude group, keep the similitude in
$1+q\mathbb Z_q$; the character is trivial there.  A sufficiently
small principal level away from $q$ is retained throughout.
It may be imposed by pullback at primes already in $S$, so
no new rational prime is introduced.  Such pullback preserves
the unit coefficient of the Eisenstein family: the same
Fourier--Jacobi value is evaluated at a lift of its cusp and
level structure.

\begin{lemma}[AC-G67-local: the deep and Iwahori sections have the same integral span]
\label{AC-G67-local}
For $m\ge1$, let $N_m$ be the upper unipotent matrices all of
whose strict upper entries lie in $q^m\mathbb Z_q$.  In a
principal series with inducing character $\alpha$, let $f_m$
be the section supported on $BwN_m$ and normalized by
$f_m(wn)=1$ for $n\in N_m$.  Assume that the restrictions of
$\alpha$ to compact diagonal units are tame, and that its
values on diagonal powers of $q$ are units in the coefficient
ring.  Then $f_m$ and $f_1$ generate the same module under
finite integral sums of right translations by
$\mathrm{GL}_3(\mathbb Q_q)$.

Hsieh Proposition~6.3 at the added $q$ gives
$\phi^{\rm pb}=\operatorname{vol}(D_W(d_qn_q))f_m$,
where $q^m=n_q$, with right diagonal character (G67.2).
This original scalar is retained in both directions of the
translation identities; no division by that volume is used.
Consequently the primitive Eisenstein family of
AC-G67-measure has a primitive integral replacement with
right $(I^-,\lambda)$-character and the same eigenvalues at
all places outside the original set $S\cup S_p$.
\end{lemma}
\begin{proof}
Put $R(u)f(g)=f(gu)$.  Since $N_m$ is a normal subgroup of
$N_1$, summing over representatives of $N_m\backslash N_1$
gives the exact identity
\[
             f_1=\sum_{u\in N_m\backslash N_1}R(u)f_m.
                                                        \tag{G67.3}
\]
To check it, restrict to the big cell $BwN$.  The unique
upper unipotent coordinate lies in $N_1$ precisely when
exactly one summand has its coordinate in $N_m$; outside
$BwN_1$ every summand vanishes.  Left $B$-equivariance gives
the identity everywhere.  No division by the index occurs.

For the reverse identity, set
\[
 t_m=\operatorname{diag}(q^{2(m-1)},q^{m-1},1),\qquad
 c_m=\alpha(wt_mw^{-1}).                                \tag{G67.4}
\]
In a normalized induction the additional modulus is a power
of $q$, also a unit here, and is included in $c_m$.
For $n=(n_{ij})\in N$,
\[
 wnt_m=(wt_mw^{-1})w(t_m^{-1}nt_m).
\]
Hence $c_m^{-1}R(t_m)f_1$ is the characteristic section of
the upper unipotent set
\[
 n_{12},n_{23}\in q^m\mathbb Z_q,\qquad
 n_{13}\in q^{2m-1}\mathbb Z_q.                         \tag{G67.5}
\]
Its nonzero value on $wN$ is $1$.  The central upper root
$u_{13}(a)=1+aE_{13}$ changes only $n_{13}$.  Therefore
\[
 f_m=c_m^{-1}
       \sum_{a\in q^m\mathbb Z_q/q^{2m-1}\mathbb Z_q}
                   R(u_{13}(a))R(t_m)f_1.                 \tag{G67.6}
\]
As in (G67.3), there is exactly one nonzero summand on the
support.  The hypothesis on $\alpha$ makes every coefficient
in (G67.3) and (G67.6) integral.

For the source identification, at $q$ the different is a
unit and the group $D_V^1(n_q)$ in (5.14) has upper entries
integral and lower entries in $q^m$.  Its diagonal endpoint
entries are $1$ modulo $q^m$; its middle diagonal entry is a
unit.  The factorization into upper unipotent, diagonal and
lower unipotent factors shows that
\[
     B w D_V(n_q)=B D_V^1(n_q) w=B w N_m.
\]
The normalization in Proposition~6.3 is $1$ on that last
unipotent set.  On compact torus elements the inducing
character is $(\vartheta_q,1,\vartheta_q^{-1})$; conjugation
by $w$ gives (G67.2).  The characters coming from the
original branch and $\Gamma_K$ are unramified at $q$, so they
do not alter this compact character.

Apply (G67.3) to the global Eisenstein family and denote the
result by $E_I$.  The ordinary projector at $p$ commutes with
these $q$-translations.  Their integrality follows either
from the displayed identities or from the prime-to-$p$
isogeny action: in Hsieh \S3.7.2 and Definition~3.13, pullback
on invariant differentials is an isomorphism because $q$ is
invertible in the coefficient ring.  Formula (G67.6) recovers
the original family from $E_I$.  If every integral coefficient
of $E_I$ lay in the coefficient maximal ideal, the same would
hold for the original family, contradicting AC-G67-measure.
For this test one may pull every translate to a common finer
level; the original unit coefficient remains a unit under
that pullback.  In particular $E_I$ is not in the maximal
ideal times the ambient ordinary weight module at $I_1^-$.
That ambient module is finite free over $\Lambda'$ by
Hsieh Theorem~4.21 (including its full, noncuspidal case) and
its weight specialization.  A coordinate
functional for a basis therefore takes a unit value on $E_I$;
its restriction to the character submodule is the integral
functional used below.  No extension of a Fourier--Jacobi
functional to a different level is being presumed.
All operators are supported at $q$, so the stated eigenvalues
away from $S\cup S_p$ are unchanged.  The normalized global
series is not renormalized in this argument: $E_I$ is defined
by the actual finite translation sum, so no factor $q-1$ is
introduced.
\end{proof}

\begin{lemma}[AC-G67-boundary: an integral character-preserving boundary lift]
\label{AC-G67-boundary}
In the imaginary quadratic Hsieh setting with $p>3$, the
ordinary boundary map at the preceding $q$-Iwahori level is
surjective on its integral $\lambda$-character submodule.
This includes the coefficient extension containing the
$p$-power values of $\lambda$.  No division by a $p$-power
group order is needed.
\end{lemma}
\begin{proof}
\proofstep{The finite deck group and its boundary stabilizers.}
Let $I_1^-$ be the kernel of the diagonal reduction
$I^-\to(\mathbb F_q^\times)^3$.  Work first at the level
with $q$-component $I_1^-$ and similitude in
$1+q\mathbb Z_q$.  The finite abelian group
\[
          A=(\mathbb F_q^\times)^3=C\times A'
                                                        \tag{G67.7}
\]
acts by changing this level structure; $C$ is its $p$-primary
part and $|A'|$ is prime to $p$.  This action commutes with
the $p$-Igusa torus, the ordinary projector and the boundary
map.  These compact-unit operators only relabel the
prime-to-$p$ level structure and preserve the Hodge bundle;
they therefore act on the minimal compactification defined
as its section-ring $\operatorname{Proj}$ in Hsieh \S2.7.3.
Retain the neat level away from $q$,
so all the schemes in question are the source's fine level
schemes.

The action of $C$ on the geometric boundary is free.  Here is
the stabilizer calculation, including the global condition.
Hsieh \S2.7.1 and Theorem~2.4(4), followed by the complex
uniformization of each definite $GU(1)$ boundary stratum,
describe its geometric points by parabolic double cosets.
If $c\in C$ fixes such a point, there are a rational Levi
element $l$, an adelic unipotent element $u$, and a compact
level element $k$ such that
\[
                   l u g c=gk.                           \tag{G67.8}
\]
The three eigenvalues of $lu$ are the rational Levi
eigenvalues $a,h,\nu\bar a^{-1}\in K^\times$.  At every
finite place, $g^{-1}lug=kc^{-1}$ lies in a compact subgroup.
Its characteristic polynomial is therefore integral and its
determinant a unit.  Each of these three elements of $K$ is
thus a unit at every finite place.  Because $K$ is imaginary
quadratic, they belong to the finite group
$\mathcal O_K^\times$, whose order divides $6$ or is $4$.
In particular all their orders are prime to $p$.

At $q$, the matrix $kc^{-1}$ is lower triangular modulo $q$.
Its three diagonal reductions are its eigenvalues modulo
$q$, so their $p$-primary components are trivial.  The
diagonal reduction of $k$ is $1$, while that of $c^{-1}$
is $c^{-1}$.  Thus $c=1$, proving freeness.

This argument applies over the coefficient base, not only
over $\mathbb C$.  By Hsieh \S2.8, the $GU(1)$ boundary
strata are finite and become disjoint copies of the
coefficient base after a finite unramified extension.  The
level maps and their deck permutations have the same
description after that extension.  Hence freeness on their
geometric points in characteristic zero proves freeness on
the geometric boundary in characteristic $p$ as well.  The
boundary quotient by $C$ is consequently a finite etale
$C$-torsor.  The same is true after restricting to the
ordinary locus.

\proofstep{A trace-one function replaces division by $|C|$.}
On a finite etale $C$-torsor, the trace of the structure
sheaf is surjective.  Indeed this can be checked after an
etale trivialization, where the algebra is a product of
$|C|$ copies and trace is the sum of the coordinates.
Because the boundary is affine, there is an integral
function $a_\partial$ there such that
\[
                    \sum_{c\in C}c(a_\partial)=1.          \tag{G67.9}
\]
Use the affine ordinary minimal space in the proof of
Hsieh Theorem~4.26.  Its affineness gives an integral lift
$a$ of $a_\partial$.  Equivalently one may work on its
$p$-adic completion.  If needed, replace the chosen lift of
the Hasse invariant by the product of all its $A$-translates;
this gives an $A$-stable affine ordinary open without
changing the ordinary formal locus.
Average $a$ over $A'$ using $|A'|^{-1}$ to make it
$A'$-invariant; (G67.9) is preserved.  The lift is a function
on the underlying ordinary space, so its pullback to the
Igusa tower has weight zero.

Let $b$ be an ordinary boundary datum of character $\lambda$.
The fundamental exact sequence supplies an integral
ordinary preimage $m$ before imposing that character.
Define on the space of $p$-adic forms
\[
 P_a(m)=\sum_{c\in C}\lambda(c)^{-1}c(a m),\qquad
 e_{A',\lambda}=|A'|^{-1}
                   \sum_{d\in A'}\lambda(d)^{-1}d.
                                                        \tag{G67.10}
\]
Changing the summation variable shows that $P_a(m)$ has
$C$-character $\lambda$.  On the boundary, since $m$ maps
to $b$, its restriction is
\[
 \sum_{c\in C}\lambda(c)^{-1}c(a_\partial)c(b)
       =\left(\sum_{c\in C}c(a_\partial)\right)b=b.
\]
The operator $e_{A',\lambda}$ preserves the $C$-character
and leaves $b$ fixed.  Finally apply the ordinary projector
$e$.  It commutes with $A$ and with the boundary operator;
since $b$ is ordinary, the boundary remains $b$.
All operations are integral and continuous.  They preserve
the prescribed Igusa weight because $a$ has weight zero.
More precisely, modulo $p^s$ the function $a$ pulls back to
$V_{n,s}=H^0(T_{n,s},\mathcal O)$ in Definition~3.4 and is
fixed by the Igusa group; multiplication therefore preserves
the weight subspace.  Section~4.3 constructs $e$ on these
complete $p$-adic modular-form modules, and Proposition~4.9
identifies their ordinary weight subspaces with ordinary
sections of the corresponding automorphic sheaf.  For every
weight in the classical dense set of (4.3), Theorem~4.19 and
Corollary~4.23 identify the resulting ordinary $p$-adic
section with a classical modular form.  The construction
commutes with reduction modulo $p^s$, with specialization of
the measure, and with the inverse limit in $s$.
It consequently satisfies the dense classical interpolation
criterion in Definition~4.24.  Theorem~4.25 now identifies
$e e_{A',\lambda}P_a(m)$ with an integral ordinary
$\Lambda'$-adic form of the required weight and character.

The initial control theorems can first be used over the
unramified coefficient ring for the fixed neat tame level.
Adjoining the values of $\lambda$ is a finite free
coefficient extension.  Tensoring with such an extension
preserves their finite freeness and exact sequences and
commutes with the displayed inverse limits and idempotents.
In particular no assumption that the tame deck group has
order prime to $p$ is introduced.  This proves the claimed
lift in the coefficient ring actually used here.
\end{proof}

\begin{proposition}[AC-G67-type: cuspidal support with finite added-prime type]
\label{AC-G67-type}
The branch in AC-G67-measure has nonzero ordinary cuspidal
Hecke support in the integral $(I^-,\lambda)$-character
module at the added prime $q$.  Every classical local
automorphic constituent in this module belongs to the tame
principal-series inertial class
\[
                (\vartheta_q^{-1},1,\vartheta_q).         \tag{G67.11}
\]
In particular its local Langlands parameter has zero
monodromy and finite inertia.  For an attached Galois
representation for which local-global compatibility at $q$
is supplied, the same finite-inertia conclusion follows.
This proposition does not assert that the resulting global
constituent is irreducible, nor any of the other away-prime
or symmetric-Hodge conditions of the actual source seed.
\end{proposition}
\begin{proof}
Put $L=L_D^S L_K^S\in\mathfrak m_{\Lambda'}$, as obtained in
AC-G67-measure.  By Hsieh Proposition~6.6, every boundary
coefficient of the original deep-level family is divisible
by $L$.  The integral translation sum (G67.3) preserves this
divisibility at every cusp, so
\[
                      \Phi(E_I)=L b                       \tag{G67.12}
\]
for an integral ordinary boundary datum $b$.
The family $E_I$ is primitive by AC-G67-local.
If $L=0$, it is already a nonzero cusp form and the assertion
follows.  Otherwise $\Lambda'$ is a domain, so (G67.12) and
the $\lambda$-transformation of $E_I$ imply the same
transformation law for $b$.

AC-G67-boundary supplies an integral ordinary
$\lambda$-character form $E_0$ with $\Phi(E_0)=b$.  Thus
\[
                       F=E_I-L E_0                       \tag{G67.13}
\]
is a cusp form of this exact character.  The integral
functional furnished in AC-G67-local is a unit on $F$, as
it is a unit on $E_I$ and $L$ lies in the maximal ideal.
Let $h_\lambda$ be the image of the Hecke algebra away from
$S\cup S_p$ on this cuspidal character module.  For every
such Hecke operator $T$, the Eisenstein relation gives
\[
       T F\equiv\lambda_E(T)F\pmod{L}.
\]
Applying the unit functional shows that the assignment
$T\mapsto\lambda_E(T)$ descends to a surjective algebra
map from the quotient of $h_\lambda$ by its Eisenstein
ideal to $\Lambda'/(L)$.  It is nonzero.  This is the exact
same algebraic subtraction argument as in Hsieh
Theorem~7.18, now inside the integral character module.

For the last assertion, use Roche, \emph{Types and Hecke
algebras for principal series representations of split
reductive $p$-adic groups}, Remark~4.2 and Theorem~7.7:
the tame torus character inflated to an Iwahori is a type
for its principal-series inertial class.  These results
apply to $\mathrm{GL}_3(\mathbb Q_q)$; the residual
characteristic here is $q$, and our choice has $q>3$.
The three compact characters in (G67.11) are pairwise
distinct, since $p>3$ and $\vartheta_q$ has order $p^r$.
Thus no two characters in that inertial class differ by
an unramified character.  More explicitly, the root system
and Weyl stabilizer attached to this torus character in
Roche \S8 are trivial; the associated group is the torus.
His Theorem~9.4 identifies parabolic induction with induction
for that torus, so these principal series are irreducible.
The associated standard principal-series
parameter is a sum of three characters with these distinct
inertial restrictions.  Its nilpotent monodromy commutes
with inertia, hence is diagonal on these three distinct
one-dimensional eigenspaces and therefore is zero.
This also gives finite actual inertia after any stated
local-global compatibility input: the unipotent factor
$\exp(t_p(\cdot)N)$ is then the identity.
\end{proof}

\paragraph{Exact remaining source-seed conditions.}
The new conclusion concerns the added split prime $q$.
After a finite abelian base change killing $\vartheta_q$,
this local parameter becomes unramified.  At $269$ the two
actual endpoint inertia characters
$\chi^j\psi$ and $\chi^j\psi^{-1}$ can coincide when
$\psi^2|_I=1$.  In that case inertia alone permits a nonzero
nilpotent map between those two endpoints.  For example,
on the common inertia eigenspace the matrix $N=E_{13}$
commutes with inertia; a diagonal Frobenius with endpoint
ratio equal to the local residue cardinality satisfies the
Weil--Deligne scaling relation.  The preceding distinct-type
argument gives no vanishing of this $N$.  Thus every
remaining away-prime \path{IsDiagOn} clause, in particular
this $269$ case, still needs its own proof.

Likewise the literal symmetric ordinary Hodge-order
obstruction in AC-G64-weight remains.  Nonzero cuspidal
support in the new finite-character branch does not by
itself establish the nonexceptional height-one-prime
hypotheses used for the earlier generic-irreducibility
argument, a source-compatible seed, or the characteristic-$p$
containment $J\subset\ker f$ in AC-G61.

\paragraph{Geometric dictionary and source scope.}
The new lifting argument uses the actual finite deck action
on the zero-dimensional boundary, not an assertion that
invariants under a $p$-group are exact.  Finite etale descent
belongs to the previously audited SGA~1, Expos\'e~V,
\S\S4 and~7 framework.  Hsieh Theorem~2.4 and \S2.8 provide
the concrete boundary geometry used here.  No additional
Grothendieck archive cote or unverified compactification
theorem is claimed.  Roche's type theorem is a retained
published input.  The Galois comparison is stated only under
the separately supplied local-global compatibility condition;
this module does not supply or verify that condition.
